% Clase 3 — Introducción a Agentes Lógicos
% Lectura: AIMA Cap. 7
\section{Agentes Lógicos: Introducción}

\subsection{Agentes Basados en Conocimiento}

Un agente basado en conocimiento mantiene una \textbf{base de conocimiento} (KB) y utiliza \textbf{inferencia} para decidir qué acción tomar.

$$\text{KB} \vdash \alpha \quad \text{(KB implica } \alpha\text{)}$$

\subsection{Lógica Proposicional}

Conectivos lógicos: $\neg$, $\land$, $\lor$, $\Rightarrow$, $\Leftrightarrow$

\begin{itemize}
    \item Una \textbf{oración} es \textit{válida} (tautología) si es verdadera en todos los modelos.
    \item Una oración es \textit{satisfacible} si es verdadera en al menos un modelo.
    \item \textbf{Consecuencia lógica}: $KB \models \alpha$ si en todo modelo donde KB es verdadera, $\alpha$ también lo es.
\end{itemize}

\subsection{Formas Normales}

\begin{itemize}
    \item \textbf{CNF} (Forma Normal Conjuntiva): conjunción de cláusulas (disyunciones de literales).
    \item \textbf{DNF} (Forma Normal Disyuntiva): disyunción de conjunciones.
\end{itemize}

% Agrega aquí el ejemplo del Mundo del Wumpus
